\Chapter{32}

\Verse{1}
{
    \zA{话说宝玉见那麒麟，心中甚是欢喜，便伸手来拿，笑道：“亏你拣着了!你是 怎么拾着的？”}
}
{
    \eA{While trying to conceal her sense of shame and injury Chin Ch'uan is driven
        by her impetuous feelings to seek death. But to resume our narrative. At the
        sight of the unicorn, Pao-yü was filled with intense delight.}
}
{
}

\Verse{2}
{
    \zA{湘云笑道：“幸而是这个。}
}
{
    \eA{So much so, that he forthwith put out his hand and made a grab for it.}
}
{
}

\Verse{3}
{
    \zA{明日倘或把印也丢了，难道也就罢了不 成？”}
}
{
    \eA{"Lucky enough it was you who picked it up!" he said, with a face beaming
        with smiles. "But when did you find it?" "Fortunately it was only this!"}
}
{
}

\Verse{4}
{
    \zA{宝玉笑道：“倒是丢了印平常，若丢了这个，我就该死了。”}
}
{
    \eA{rejoined Shih Hsiang-yün laughing. "If you by and bye also lose your seal,
        will you likely banish it at once from your mind, and never make an effort
        to discover it?"}
}
{
}

\Verse{5}
{
    \zA{袭人倒了茶来与湘云吃，一面笑道：“大姑娘，我前日听见你大喜呀。”}
}
{
    \eA{"After all," smiled Pao-yü, "the loss of a seal is an ordinary occurrence.
        But had I lost this, I would have deserved to die." Hsi Jen then poured a
        cup of tea and handed it to Shih Hsiang-yün.}
}
{
}

\Verse{6}
{
    \zA{湘云 红了脸，扭过头去吃茶，一声也不答应。}
}
{
    \eA{"Miss Senior," she remarked smilingly, "I heard that you had occasion the
        other day to be highly pleased." Shih Hsiang-yün flushed crimson. She went
        on drinking her tea and did not utter a single word.}
}
{
}

\Verse{7}
{
    \zA{袭人笑道：“这会子又害臊了？}
}
{
    \eA{"Here you are again full of shame!" Hsi Jen smiled. "But do you remember
        when we were living,}
}
{
}

\Verse{8}
{
    \zA{你还记得 那几年，咱们在西边暖阁上住着，晚上你和我说的话？}
}
{
    \eA{about ten years back, in those warm rooms on the west side and you confided
        in me one evening, you didn't feel any shame then; and how is it you blush
        like this now?"}
}
{
}

\Verse{9}
{
    \zA{那会子不害臊，这会子怎么 又臊了？”}
}
{
    \eA{"Do you still speak about that!" exclaimed Shih Hsiang-yün laughingly.}
}
{
}

\Verse{10}
{
    \zA{湘云的脸越发红了，勉强笑道：“你还说呢!那会子咱们那么好，后来 我们太太没了，我家去住了一程子，怎么就把你配给了他。}
}
{
    \eA{"You and I were then great friends. But when our mother subsequently died
        and I went home for a while, how is it you were at once sent to be with my
        cousin Secundus,}
}
{
}

\Verse{11}
{
    \zA{我来了，你就不那么待 我了。”}
}
{
    \eA{and that now that I've come back you don't treat me as you did once?" "Are
        you yet harping on this!"}
}
{
}

\Verse{12}
{
    \zA{袭人也红了脸，笑道：“罢呦!先头里，‘姐姐’长，‘姐姐’短，哄着 我替你梳头洗脸，做这个弄那个，如今拿出小姐款儿来了。}
}
{
    \eA{retorted Hsi Jen, putting on a smile. "Why, at first, you used to coax me
        with a lot of endearing terms to comb your hair and to wash your face, to do
        this and that for you. But now that you've become a big girl,}
}
{
}

\Verse{13}
{
    \zA{你既拿款，我敢亲近 吗？”}
}
{
    \eA{you assume the manner of a young mistress towards me,}
}
{
}

\Verse{14}
{
    \zA{湘云道：“阿弥陀佛，冤枉冤哉!我要这么着，就立刻死了。}
}
{
    \eA{and as you put on these airs of a young mistress, how can I ever presume to
        be on a familiar footing with you?" "O-mi-to-fu," cried Shih Hsiang-yün.}
}
{
}

\Verse{15}
{
    \zA{你瞧瞧，这么 大热天，我来了必定先瞧瞧你。}
}
{
    \eA{"What a false accusation! If I be guilty of anything of the kind, may I at
        once die! Just see what a broiling hot day this is,}
}
{
}

\Verse{16}
{
    \zA{你不信问缕儿：我在家时时刻刻，那一回不想念你 几句？”}
}
{
    \eA{and yet as soon as I arrived I felt bound to come and look you up first. If
        you don't believe me, well, ask Lü Erh! And while at home, when did I not at
        every instant say something about you?"}
}
{
}

\Verse{17}
{
    \zA{袭人和宝玉听了，都笑劝道：“说玩话儿，你又认真了。}
}
{
    \eA{Scarcely had she concluded than Hsi Jen and Pao-yü tried to soothe her. "We
        were only joking," they said, "but you've taken everything again as gospel.}
}
{
}

\Verse{18}
{
    \zA{还是这么性儿急。”}
}
{
    \eA{What! are you still so impetuous in your temperament!"}
}
{
}

\Verse{19}
{
    \zA{湘云道：“你不说你的话咽人，倒说人性急。”}
}
{
    \eA{"You don't say," argued Shih Hsiang-yün, "that your words are hard things to
        swallow, but contrariwise, call people's temperaments impetuous!"}
}
{
}

\Verse{20}
{
    \zA{一面说，一面打开绢子，将戒指递与袭人。}
}
{
    \eA{As she spoke, she unfolded her handkerchief and, producing a ring, she gave
        it to Hsi Jen.}
}
{
}

\Verse{21}
{
    \zA{袭人感谢不尽，因笑道：“你前日 送你姐姐们的，我已经得了。}
}
{
    \eA{Hsi Jen did not know how to thank her enough. "When;" she consequently
        smiled, "you sent those to your cousin the other day,}
}
{
}

\Verse{22}
{
    \zA{今日你亲自又送来，可见是没忘了我。}
}
{
    \eA{I got one also; and here you yourself bring me another to-day! It's clear
        enough therefore that you haven't forgotten me.}
}
{
}

\Verse{23}
{
    \zA{就为这个试出 你来了。}
}
{
    \eA{This alone has been quite enough to test you.}
}
{
}

\Verse{24}
{
    \zA{戒指儿能值多少，可见你的心真。”}
}
{
    \eA{As for the ring itself, what is its worth? but it's a token of the sincerity
        of your heart!"}
}
{
}

\Verse{25}
{
    \zA{史湘云道：“是谁给你的？”}
}
{
    \eA{"Who gave it to you?" inquired Shih Hsiang-yün. "Miss Pao let me have it."}
}
{
}

\Verse{26}
{
    \zA{袭人道： “是宝姑娘给我的。”}
}
{
    \eA{replied Hsi Jen. "I was under the impression," remarked Hsiang-yün with a
        sigh,}
}
{
}

\Verse{27}
{
    \zA{湘云叹道：“我只当林姐姐送你的，原来是宝姐姐给了你。}
}
{
    \eA{"that it was a present from cousin Lin. But is it really cousin Pao, that
        gave it to you! When I was at home,}
}
{
}

\Verse{28}
{
    \zA{我天天在家里想着，这些姐姐们，再没一个比宝姐姐好的。}
}
{
    \eA{I day after day found myself reflecting that among all these cousins of
        mine, there wasn't one able to compare with cousin Pao,}
}
{
}

\Verse{29}
{
    \zA{可惜我们不是一个娘养 的。}
}
{
    \eA{so excellent is she. How I do regret that we are not the offspring of one
        mother!}
}
{
}

\Verse{30}
{
    \zA{我但凡有这么个亲姐姐，就是没了父母，也没妨碍的！”}
}
{
    \eA{For could I boast of such a sister of the same flesh and blood as myself, it
        wouldn't matter though I had lost both father and mother!"}
}
{
}

\Verse{31}
{
    \zA{说道，眼圈儿就红了。}
}
{
    \eA{While indulging in these regrets, her eyes got quite red. "Never mind!}
}
{
}

\Verse{32}
{
    \zA{宝玉道：“罢罢罢，不用提起这个话了。”}
}
{
    \eA{never mind!" interposed Pao-yü. "Why need you speak of these things!" "If I
        do allude to this,"}
}
{
}

\Verse{33}
{
    \zA{史湘云道：“提这个便怎么？}
}
{
    \eA{answered Shih Hsiang-yün, "what does it matter? I know that weak point of
        yours.}
}
{
}

\Verse{34}
{
    \zA{我知道你 的心病：恐怕你的林妹妹听见，又嗔我赞了宝姐姐了。}
}
{
    \eA{You're in fear and trembling lest your cousin Lin should come to hear what I
        say, and get angry with me again for eulogising cousin Pao!}
}
{
}

\Verse{35}
{
    \zA{可是为这个不是？”}
}
{
    \eA{Now isn't it this, eh!" "Ch'ih!"}
}
{
}

\Verse{36}
{
    \zA{袭人在 旁嗤的一笑，说道：“云姑娘，你如今大了，越发心直嘴快了。”}
}
{
    \eA{laughed Hsi Jen, who was standing by her. "Miss Yün," she said, "now that
        you've grown up to be a big girl you've become more than ever openhearted
        and outspoken."}
}
{
}

\Verse{37}
{
    \zA{宝玉笑道：“我 说你们这几个人难说话，果然不错。”}
}
{
    \eA{"When I contend;" smiled Pao-yü, "that it is difficult to say a word to any
        one of you I'm indeed perfectly correct!"}
}
{
}

\Verse{38}
{
    \zA{史湘云道：“好哥哥，你不必说话叫我恶心。}
}
{
    \eA{"My dear cousin," observed Shih Hsiang-yün laughingly, "don't go on in that
        strain!}
}
{
}

\Verse{39}
{
    \zA{只会在我跟前说话，见了你林妹妹，又不知怎么好了。”}
}
{
    \eA{You'll provoke me to displeasure. When you are with me all you are good for
        is to talk and talk away; but were you to catch a glimpse of cousin Lin,}
}
{
}

\Verse{40}
{
    \zA{袭人道：“且别说玩话，正有一件事要求你呢。”}
}
{
    \eA{you would once more be quite at a loss to know what best to do!" "Now,
        enough of your jokes!" urged Hsi Jen. "I have a favour to crave of you."}
}
{
}

\Verse{41}
{
    \zA{史湘云便问：“什么事？”}
}
{
    \eA{"What is it?" vehemently inquired Shih Hsiang-yün.}
}
{
}

\Verse{42}
{
    \zA{袭人道：“有一双鞋，抠了垫心子，我这两日身上不好，不得做，你可有工夫替我 做做？”}
}
{
    \eA{"I've got a pair of shoes," answered Hsi Jen, "for which I've stuck the
        padding together; but I'm not feeling up to the mark these last few days, so
        I haven't been able to work at them.}
}
{
}

\Verse{43}
{
    \zA{史湘云道：“这又奇了。}
}
{
    \eA{If you have any leisure, do finish them for me." "This is indeed strange!"}
}
{
}

\Verse{44}
{
    \zA{你家放着这些巧人不算，还有什么针线上的、裁 剪上的，怎么叫我做起来？}
}
{
    \eA{exclaimed Shih Hsiang-yün. "Putting aside all the skilful workers engaged in
        your household, you have besides some people for doing needlework and others
        for tailoring and cutting;}
}
{
}

\Verse{45}
{
    \zA{你的活计叫人做，谁好意思不做呢？”}
}
{
    \eA{and how is it you appeal to me to take your shoes in hand? Were you to ask
        any one of those men to execute your work,}
}
{
}

\Verse{46}
{
    \zA{袭人笑道：“你 又糊涂了。}
}
{
    \eA{who could very well refuse to do it?" "Here you are in another stupid mood!"}
}
{
}

\Verse{47}
{
    \zA{你难道不知道：我们这屋里的针线，是不要那些针线上的人做的。”}
}
{
    \eA{laughed Hsi Jen. "Can it be that you don't know that our sewing in these
        quarters mayn't be done by these needleworkers."}
}
{
}

\Verse{48}
{
    \zA{史 湘云听了，便知是宝玉的鞋，因笑道：“既这么说，我就替你做做罢。}
}
{
    \eA{At this reply, it at once dawned upon Shih Hsiang-yün that the shoes must be
        intended for Pao-yü. "Since that be the case," she in consequence smiled;
        "I'll work them for you.}
}
{
}

\Verse{49}
{
    \zA{只是一件： 你的我才做，别人的我可不能。”}
}
{
    \eA{There's however one thing. I'll readily attend to any of yours, but I will
        have nothing to do with any for other people."}
}
{
}

\Verse{50}
{
    \zA{袭人笑道：“又来了。}
}
{
    \eA{"There you are again!" laughed Hsi Jen.}
}
{
}

\Verse{51}
{
    \zA{我是个什么儿，就敢烦你 做鞋了!实告诉你：可不是我的。}
}
{
    \eA{"Who am I to venture to trouble you to make shoes for me? I'll tell you
        plainly, however, that they are not mine. But no matter whose they are,}
}
{
}

\Verse{52}
{
    \zA{你别管是谁的，横竖我领情就是了。”}
}
{
    \eA{it is anyhow I who'll be the recipient of your favour; that is sufficient."
        "To speak the truth,"}
}
{
}

\Verse{53}
{
    \zA{史湘云道： “论理，你的东西也不知烦我做了多少。}
}
{
    \eA{rejoined Shih Hsiang-yün, "you've put me to the trouble of working, I don't
        know how many things for you.}
}
{
}

\Verse{54}
{
    \zA{今日我倒不做的原故，你必定也知道。”}
}
{
    \eA{The reason why I refuse on this occasion should be quite evident to you!" "I
        can't nevertheless make it out!"}
}
{
}

\Verse{55}
{
    \zA{袭人道：“我倒也不知道。”}
}
{
    \eA{answered Hsi Jen. "I heard the other day," continued Shih Hsiang-yün,}
}
{
}

\Verse{56}
{
    \zA{史湘云冷笑道：“前日我听见把我做的扇套儿拿着和 人家比，赌气又铰了。}
}
{
    \eA{a sardonic smile on her lip, "that while the fan-case, I had worked, was
        being held and compared with that of some one else, it too was slashed away
        in a fit of high dudgeon. This reached my ears long ago,}
}
{
}

\Verse{57}
{
    \zA{我早就听见了，你还瞒我？}
}
{
    \eA{and do you still try to dupe me by asking me again now to make something
        more for you?}
}
{
}

\Verse{58}
{
    \zA{这会子又叫我做，我成了你们奴 才了。”}
}
{
    \eA{Have I really become a slave to you people? "As to what occurred the other
        day," hastily explained Pao-yü smiling,}
}
{
}

\Verse{59}
{
    \zA{宝玉忙笑道：“前日的那个本不知是你做的。”}
}
{
    \eA{"I positively had no idea that that thing was your handiwork." "He never
        knew that you'd done it,"}
}
{
}

\Verse{60}
{
    \zA{袭人也笑道：“他本不知 是你做的，是我哄他的话，说是‘新近外头有个会做活的，扎的绝出奇的好花儿，
        叫他们拿了一个扇套儿试试看好不好’，他就信了，拿出去给这个瞧、那个看的。}
}
{
    \eA{Hsi Jen also laughed. "I deceived him by telling him that there had been of
        late some capital hands at needlework outside, who could execute any
        embroidery with surpassing beauty, and that I had asked them to bring a
        fan-case so as to try them and to see whether they could actually work well
        or not. He at once believed what I said. But as he produced the case and
        gave it to this one and that one to look at,}
}
{
}

\Verse{61}
{
    \zA{不知怎么又惹恼了那一位，铰了两段。}
}
{
    \eA{he somehow or other, I don't know how, managed again to put some one's back
        up, and she cut it into two.}
}
{
}

\Verse{62}
{
    \zA{回来他还叫赶着做去，我才说了是你做的， 他后悔的什么似的！”}
}
{
    \eA{On his return, however, he bade me hurry the men to make another; and when
        at length I explained to him that it had been worked by you, he felt, I
        can't tell you,}
}
{
}

\Verse{63}
{
    \zA{史湘云道：“这越发奇了。}
}
{
    \eA{what keen regret!" "This is getting stranger and stranger!"}
}
{
}

\Verse{64}
{
    \zA{林姑娘也犯不上生气，他既会剪， 就叫他做。”}
}
{
    \eA{said Shih Hsiang-yün. "It wasn't worth the while for Miss Lin to lose her
        temper about it. But as she plies the scissors so admirably,}
}
{
}

\Verse{65}
{
    \zA{袭人道：“他可不做呢。}
}
{
    \eA{why, you might as well tell her to finish the shoes for you." "She
        couldn't," replied Hsi Jen,}
}
{
}

\Verse{66}
{
    \zA{饶这么着老太太还怕他劳碌着了，大夫又说 好生静养才好，谁还肯烦他做呢？}
}
{
    \eA{"for besides other things our venerable lady is still in fear and trembling
        lest she should tire herself in any way. The doctor likewise says that she
        will continue to enjoy good health, so long as she is carefully looked
        after; so who would wish to ask her to take them in hand? Last year she
        managed to just get through a scented bag,}
}
{
}

\Verse{67}
{
    \zA{旧年好一年的工夫做了个香袋儿，今年半年还没 见拿针线呢。”}
}
{
    \eA{after a whole year's work. But here we've already reached the middle of the
        present year, and she hasn't yet taken up any needle or thread!" In the
        course of their conversation,}
}
{
}

\Verse{68}
{
    \zA{正说着，有人来回说：“兴隆街的大爷来了，老爷叫二爷出去会。”}
}
{
    \eA{a servant came and announced 'that the gentleman who lived in the Hsing Lung
        Street had come.' "Our master," he added, "bids you, Mr. Secundus, come out
        and greet him."}
}
{
}

\Verse{69}
{
    \zA{宝玉听了， 便知贾雨村来了，心中好不自在。}
}
{
    \eA{As soon as Pao-yü heard this announcement, he knew that Chia Yü-ts'un must
        have arrived. But he felt very unhappy at heart.}
}
{
}

\Verse{70}
{
    \zA{袭人忙去拿衣服。}
}
{
    \eA{Hsi Jen hurried to go and bring his clothes.}
}
{
}

\Verse{71}
{
    \zA{宝玉一面登着靴子，一面抱怨 道：“有老爷和他坐着就罢了，回回定要见我！”}
}
{
    \eA{Pao-yü, meanwhile, put on his boots, but as he did so, he gave way to
        resentment. "Why there's father," he soliloquised, "to sit with him; that
        should be enough; and must he, on every visit he pays,}
}
{
}

\Verse{72}
{
    \zA{史湘云一边摇着扇子，笑道：“自 然你能迎宾接客，老爷才叫你出去呢。”}
}
{
    \eA{insist upon seeing me!" "It is, of course, because you have such a knack for
        receiving and entertaining visitors that Mr. Chia Cheng will have you go
        out," laughingly interposed Shih Hsiang-yün from one side, as she waved her
        fan. "Is it father's doing?"}
}
{
}

\Verse{73}
{
    \zA{宝玉道：“那里是老爷？}
}
{
    \eA{Pao-yü rejoined. "Why, it's he himself who asks that I should be sent for to
        see him."}
}
{
}

\Verse{74}
{
    \zA{都是他自己要请 我见的。”}
}
{
    \eA{"'When a host is courteous, visitors come often,'" smiled Hsiang-yün,}
}
{
}

\Verse{75}
{
    \zA{湘云笑道：“‘主雅客来勤’，自然你有些警动他的好处，他才要会你。”}
}
{
    \eA{"so it's surely because you possess certain qualities, which have won his
        regard, that he insists upon seeing you." "But I am not what one would call
        courteous," demurred Pao-yü. "I am, of all coarse people,}
}
{
}

\Verse{76}
{
    \zA{宝玉道：“罢，罢，我也不过俗中又俗的一个俗人罢了，并不愿和这些人来往。”}
}
{
    \eA{the coarsest. Besides, I do not choose to have any relations with such
        people as himself." "Here's again that unchangeable temperament of yours!"
        laughed Hsiang-yün. "But you're a big fellow now, and you should at least,}
}
{
}

\Verse{77}
{
    \zA{湘云笑道：“还是这个性儿，改不了!如今大了，你就不愿意去考举人进士的，也 该常会会这些为官作宦的，谈讲谈讲那些仕途经济，也好将来应酬事务，日后也有
        个正经朋友。}
}
{
    \eA{if you be loth to study and go and pass your examinations for a provincial
        graduate or a metropolitan graduate, have frequent intercourse with officers
        and ministers of state and discuss those varied attainments, which one
        acquires in an official career, so that you also may be able in time to have
        some idea about matters in general; and that when by and bye you've made
        friends, they may not see you spending the whole day long in doing nothing
        than loafing in our midst,}
}
{
}

\Verse{78}
{
    \zA{让你成年家只在我们队里，搅的出些什么来？”}
}
{
    \eA{up to every imaginable mischief." "Miss," exclaimed Pao-yü, after this
        harangue, "pray go and sit in some other girl's room,}
}
{
}

\Verse{79}
{
    \zA{宝玉听了，大觉逆耳，便道：“姑娘请别的屋里坐坐罢，我这里仔细腌？}
}
{
    \eA{for mind one like myself may contaminate a person who knows so much of
        attainments and experience as you do." "Miss," ventured Hsi Jen, "drop this
        at once! Last time Miss Pao too tendered him this advice, but without
        troubling himself as to whether people would feel uneasy or not,}
}
{
}

\Verse{80}
{
    \zA{了你 这样知经济的人！”}
}
{
    \eA{he simply came out with an ejaculation of 'hai,' and rushed out of the
        place.}
}
{
}

\Verse{81}
{
    \zA{袭人连忙解说道：“姑娘快别说他。}
}
{
    \eA{Miss Pao hadn't meanwhile concluded her say, so when she saw him fly, she
        got so full of shame that,}
}
{
}

\Verse{82}
{
    \zA{上回也是宝姑娘说过一回， 他也不管人脸上过不去，？}
}
{
    \eA{flushing scarlet, she could neither open her lips, nor hold her own counsel.
        But lucky for him it was only Miss Pao.}
}
{
}

\Verse{83}
{
    \zA{了一声，拿起脚来就走了。}
}
{
    \eA{Had it been Miss Lin, there's no saying what row there may not have been
        again, and what tears may not have been shed!}
}
{
}

\Verse{84}
{
    \zA{宝姑娘的话也没说完，见他 走了，登时羞的脸通红，说不是，不说又不是。}
}
{
    \eA{Yet the very mention of all she had to tell him is enough to make people
        look up to Miss Pao with respect. But after a time, she also betook herself
        away. I then felt very unhappy as I imagined that she was angry; but
        contrary to all my expectations,}
}
{
}

\Verse{85}
{
    \zA{幸而是宝姑娘，那要是林姑娘，不 知又闹的怎么样、哭的怎么样呢!提起这些话来，宝姑娘叫人敬重。}
}
{
    \eA{she was by and bye just the same as ever. She is, in very truth,
        long-suffering and indulgent! This other party contrariwise became quite
        distant to her, little though one would have thought it of him; and as Miss
        Pao perceived that he had lost his temper,}
}
{
}

\Verse{86}
{
    \zA{自己过了一会 子去了，我倒过不去，只当他恼了，谁知过后还是照旧一样，真真是有涵养、心地 宽大的。}
}
{
    \eA{and didn't choose to heed her, she subsequently made I don't know how many
        apologies to him." "Did Miss Lin ever talk such trash!" exclaimed Pao-yü.
        "Had she ever talked such stuff and nonsense, I would have long ago become
        chilled towards her."}
}
{
}

\Verse{87}
{
    \zA{谁知这一位反倒和他生分了。}
}
{
    \eA{"What you say is all trash!" Hsi Jen and Hsiang-yün remarked with one voice,}
}
{
}

\Verse{88}
{
    \zA{那林姑娘见他赌气不理，他后来不知赔多少 不是呢。”}
}
{
    \eA{while they shook their heads to and fro and smiled. Lin Tai-yü, the fact is,
        was well aware that now that Shih Hsiang-yün was staying in the mansion,}
}
{
}

\Verse{89}
{
    \zA{宝玉道：“林姑娘从来说过这些混帐话吗？}
}
{
    \eA{Pao-yü too was certain to hasten to come and tell her all about the unicorn
        he had got, so she thought to herself:}
}
{
}

\Verse{90}
{
    \zA{要是他也说过这些混帐话， 我早和他生分了。”}
}
{
    \eA{"In the foreign traditions and wild stories, introduced here of late by
        Pao-yü, literary persons and pretty girls are,}
}
{
}

\Verse{91}
{
    \zA{袭人和湘云都点头笑道：“这原是混帐话么？”}
}
{
    \eA{for the most part, brought together in marriage, through the agency of some
        trifling but ingenious nick-nack. These people either have miniature ducks,}
}
{
}

\Verse{92}
{
    \zA{原来黛玉知道史湘云在这里，宝玉一定又赶来，说麒麟的原故。}
}
{
    \eA{or phoenixes, jade necklets or gold pendants, fine handkerchiefs or elegant
        sashes; and they have, through the instrumentality of such trivial objects,}
}
{
}

\Verse{93}
{
    \zA{因心下忖度着， 近日宝玉弄来的外传野史，多半才子佳人，都因小巧玩物上撮合，或有鸳鸯，或有 凤凰，或玉环金佩，或鲛帕鸾绦，皆由小物而遂终身之愿。}
}
{
    \eA{invariably succeeded in accomplishing the wishes they entertained throughout
        their lives." When she recently discovered, by some unforeseen way, that
        Pao-yü had likewise a unicorn she began to apprehend lest he should make
        this circumstance a pretext to create an estrangement with her, and indulge
        with Shih Hsiang-yün as well in various free and easy flirtations and fine
        doings. She therefore quietly crossed over to watch her opportunity and take
        such action as would enable her to get an insight into his and her
        sentiments.}
}
{
}

\Verse{94}
{
    \zA{今忽见宝玉也有麒麟， 便恐借此生隙，同湘云也做出那些风流佳事来。}
}
{
    \eA{Contrary, however, to all her calculations, no sooner did she reach her
        destination, than she overheard Shih Hsiang-yün dilate on the topic of
        experience, and Pao-yü go on to observe:}
}
{
}

\Verse{95}
{
    \zA{因而悄悄走来，见机行事，以察二 人之意。}
}
{
    \eA{"Cousin Lin has never indulged in such stuff and nonsense. Had she ever
        uttered any such trash, I would have become chilled even towards her!"}
}
{
}

\Verse{96}
{
    \zA{不想刚走进来，正听见湘云说“经济”一事，宝玉又说“林妹妹不说这些 混帐话，要说这话，我也和他生分了”。}
}
{
    \eA{This language suddenly produced, in Lin Tai-yü's mind, both surprise as well
        as delight; sadness as well as regret. Delight, at having indeed been so
        correct in her perception that he whom she had ever considered in the light
        of a true friend had actually turned out to be a true friend.}
}
{
}

\Verse{97}
{
    \zA{黛玉听了这话，不觉又喜又惊，又悲又叹。}
}
{
    \eA{Surprise, "because," she said to herself: "he has, in the presence of so
        many witnesses, displayed such partiality as to speak in my praise,}
}
{
}

\Verse{98}
{
    \zA{所喜者：果然自己眼力不错，素日认他是个知己，果然是个知己；所惊者：他在人 前一片私心称扬于我，其亲热厚密，竟不避嫌疑；所叹者：你既为我的知己，自然
        我亦可为你的知己，既你我为知己，又何必有“金玉”之论呢？}
}
{
    \eA{and has shown such affection and friendliness for me as to make no attempt
        whatever to shirk suspicion." Regret, "for since," (she pondered), "you are
        my intimate friend, you could certainly well look upon me too as your
        intimate friend; and if you and I be real friends,}
}
{
}

\Verse{99}
{
    \zA{既有“金玉”之论， 也该你我有之，又何必来一宝钗呢？}
}
{
    \eA{why need there be any more talk about gold and jade? But since there be that
        question of gold and jade, you and I should have such things in our
        possession.}
}
{
}

\Verse{100}
{
    \zA{所悲者：父母早逝，虽有铭心刻骨之言，无人 为我主张；况近日每觉神思恍惚，病已渐成，医者更云：“气弱血亏，恐致劳怯之 症。”}
}
{
    \eA{Yet, why should this Pao-ch'ai step in again between us?" Sad, "because,"
        (she reflected), "my father and mother departed life at an early period; and
        because I have, in spite of the secret engraven on my heart and imprinted on
        my bones, not a soul to act as a mentor to me. Besides, of late, I
        continuously feel confusion creep over my mind,}
}
{
}

\Verse{101}
{
    \zA{我虽为你的知己，但恐不能久待；你纵为我的知己，奈我薄命何!想到此间， 不禁泪又下来。}
}
{
    \eA{so my disease must already have gradually developed itself. The doctors
        further state that my breath is weak and my blood poor, and that they dread
        lest consumption should declare itself, so despite that sincere friendship I
        foster for you,}
}
{
}

\Verse{102}
{
    \zA{待要进去相见，自觉无味，便一面拭泪，一面抽身回去了。}
}
{
    \eA{I cannot, I fear, last for very long. You are, I admit, a true friend to me,
        but what can you do for my unfortunate destiny!" Upon reaching this point in
        her reflections,}
}
{
}

\Verse{103}
{
    \zA{这里宝玉忙忙的穿了衣裳出来，忽见黛玉在前面慢慢的走着，似乎有拭泪之 状，便忙赶着上来笑道：“妹妹往那里去？}
}
{
    \eA{she could not control her tears, and they rolled freely down her cheeks. So
        much so, that when about to enter and meet her cousins, she experienced such
        utter lack of zest, that, while drying her tears she turned round, and
        wended her steps back in the direction of her apartments.}
}
{
}

\Verse{104}
{
    \zA{怎么又哭了？}
}
{
    \eA{Pao-yü, meanwhile, had hurriedly got into his new costume.}
}
{
}

\Verse{105}
{
    \zA{又是谁得罪了你了？”}
}
{
    \eA{Upon coming out of doors, he caught sight of Lin Tai-yü,}
}
{
}

\Verse{106}
{
    \zA{黛玉 回头见是宝玉，便勉强笑道：“好好的，我何曾哭来。”}
}
{
    \eA{walking quietly ahead of him engaged, to all appearances, in wiping tears
        from her eyes. With rapid stride, he overtook her. "Cousin Lin," he smiled,}
}
{
}

\Verse{107}
{
    \zA{宝玉笑道：“你瞧瞧，睛 睛上的泪珠儿没干，还撒谎呢。”}
}
{
    \eA{"where are you off to? How is it that you're crying again? Who has once more
        hurt your feelings?" Lin Tai-yü turned her head round to look;}
}
{
}

\Verse{108}
{
    \zA{一面说，一面禁不住抬起手来，替他拭泪。}
}
{
    \eA{and seeing that it was Pao-yü, she at once forced a smile. "Why should I be
        crying,"}
}
{
}

\Verse{109}
{
    \zA{黛玉 忙向后退了几步，说道：“你又要死了!又这么动手动脚的。”}
}
{
    \eA{she replied, "when there is no reason to do so?" "Look here!" observed
        Pao-yü smilingly. "The tears in your eyes are not dry yet and do you still
        tell me a fib?"}
}
{
}

\Verse{110}
{
    \zA{宝玉笑道：“说话 忘了情，不觉的动了手，也就顾不得死活。”}
}
{
    \eA{Saying this, he could not check an impulse to raise his arm and wipe her
        eyes, but Lin Tai-yü speedily withdrew several steps backwards. "Are you
        again bent," she said,}
}
{
}

\Verse{111}
{
    \zA{黛玉道：“死了倒不值什么，只是丢 下了什么‘金’，又是什么‘麒麟’，可怎么好呢！”}
}
{
    \eA{"upon compassing your own death! Then why do you knock your hands and kick
        your feet about in this wise?" "While intent upon speaking, I forgot,"
        smiled Pao-yü, "all about propriety and gesticulated,}
}
{
}

\Verse{112}
{
    \zA{一句话又把宝玉说急了，赶 上来问道：“你还说这些话，到底是咒我还是气我呢？”}
}
{
    \eA{yet quite inadvertently. But what care I whether I die or live!" "To die
        would, after all" added Lin Tai-yü, "be for you of no matter; but you'll
        leave behind some gold or other, and a unicorn too or other;}
}
{
}

\Verse{113}
{
    \zA{黛玉见问，方想起前日的 事来，遂自悔这话又说造次了，忙笑道：“你别着急，我原说错了。}
}
{
    \eA{and what would they do?" This insinuation was enough to plunge Pao-yü into a
        fresh fit of exasperation. Hastening up to her: "Do you still give vent to
        such language?" he asked. "Why, it's really tantamount to invoking
        imprecations on me!}
}
{
}

\Verse{114}
{
    \zA{这有什么要紧， 筋都叠暴起来，急的一脸汗!”一面说，一面也近前伸手替他拭面上的汗。}
}
{
    \eA{What, are you yet angry with me!" This question recalled to Lin Tai-yü's
        mind the incidents of a few days back, and a pang of remorse immediately
        gnawed her heart for having been again so indiscreet in her speech.}
}
{
}

\Verse{115}
{
    \zA{宝玉瞅了半天，方说道：“你放心。”}
}
{
    \eA{"Now don't you distress your mind!" she observed hastily, smiling. "I verily
        said what I shouldn't!}
}
{
}

\Verse{116}
{
    \zA{黛玉听了，怔了半天，说道：“我有什 么不放心的？}
}
{
    \eA{Yet what is there in this to make your veins protrude, and to so provoke you
        as to bedew your whole face with perspiration?"}
}
{
}

\Verse{117}
{
    \zA{我不明白你这个话。}
}
{
    \eA{While reasoning with him, she felt unable to repress herself,}
}
{
}

\Verse{118}
{
    \zA{你倒说说，怎么放心不放心？”}
}
{
    \eA{and, approaching him, she extended her hand, and wiped the perspiration from
        his face.}
}
{
}

\Verse{119}
{
    \zA{宝玉叹了一口气， 问道：“你果然不明白这话？}
}
{
    \eA{Pao-yü gazed intently at her for a long time. "Do set your mind at ease!" he
        at length observed. At this remark,}
}
{
}

\Verse{120}
{
    \zA{难道我素日在你身上的心都用错了？}
}
{
    \eA{Lin Tai-yü felt quite nervous. "What's there to make my mind uneasy?" she
        asked after a protracted interval.}
}
{
}

\Verse{121}
{
    \zA{连你的意思若体贴 不着，就难怪你天天为我生气了。”}
}
{
    \eA{"I can't make out what you're driving at; tell me what's this about making
        me easy or uneasy?" Pao-yü heaved a sigh. "Don't you truly fathom the depth
        of my words?"}
}
{
}

\Verse{122}
{
    \zA{黛玉道：“我真不明白放心不放心的话。”}
}
{
    \eA{he inquired. "Why, do you mean to say that I've throughout made such poor
        use of my love for you as not to be able to even divine your feelings?}
}
{
}

\Verse{123}
{
    \zA{宝 玉点头叹道：“好妹妹，你别哄我。}
}
{
    \eA{Well, if so, it's no wonder that you daily lose your temper on my account!"}
}
{
}

\Verse{124}
{
    \zA{你真不明白这话，不但我素日白用了心，且连 你素日待我的心也都辜负了。}
}
{
    \eA{"I actually don't understand what you mean by easy or uneasy," Lin Tai-yü
        replied. "My dear girl," urged Pao-yü, nodding and sighing. "Don't be making
        a fool of me! For if you can't make out these words,}
}
{
}

\Verse{125}
{
    \zA{你皆因都是不放心的原故，才弄了一身的病了。}
}
{
    \eA{not only have I ever uselessly lavished affection upon you, but the regard,
        with which you have always treated me, has likewise been entirely of no
        avail!}
}
{
}

\Verse{126}
{
    \zA{但凡 宽慰些，这病也不得一日重似一日了！”}
}
{
    \eA{And it's mostly because you won't set your mind at ease that your whole
        frame is riddled with disease. Had you taken things easier a bit,}
}
{
}

\Verse{127}
{
    \zA{黛玉听了这话，如轰雷掣电，细细思之，竟比自己肺腑中掏出来的还觉恳切， 竟有万句言语，满心要说，只是半个字也不能吐出，只管怔怔的瞅着他。}
}
{
    \eA{this ailment of yours too wouldn't have grown worse from day to day!" These
        words made Lin Tai-yü feel as if she had been blasted by thunder, or struck
        by lightning. But after carefully weighing them within herself, they seemed
        to her far more fervent than any that might have emanated from the depths of
        her own heart, and thousands of sentiments, in fact, thronged together in
        her mind;}
}
{
}

\Verse{128}
{
    \zA{此时宝玉 心中也有万句言词，不知一时从那一句说起，却也怔怔的瞅着黛玉。}
}
{
    \eA{but though she had every wish to frame them into language, she found it a
        hard task to pronounce so much as half a word. All she therefore did was to
        gaze at him with vacant stare. Pao-yü fostered innumerable thoughts within
        himself,}
}
{
}

\Verse{129}
{
    \zA{两个人怔了半 天，黛玉只？}
}
{
    \eA{but unable in a moment to resolve from which particular one to begin,}
}
{
}

\Verse{130}
{
    \zA{了一声，眼中泪直流下来，回身便走。}
}
{
    \eA{he too absently looked at Tai-yü. Thus it was that the two cousins remained
        for a long time under the spell of a deep reverie. An ejaculation of "Hai!"}
}
{
}

\Verse{131}
{
    \zA{宝玉忙上前拉住道：“好妹妹， 且略站住，我说一句话再走。”}
}
{
    \eA{was the only sound that issued from Lin Tai-yü's lips; and while tears
        streamed suddenly from her eyes, she turned herself round and started on her
        way homeward.}
}
{
}

\Verse{132}
{
    \zA{黛玉一面拭泪，一面将手推开，说道：“有什么可 说的？}
}
{
    \eA{Pao-yü jumped forward, with alacrity, and dragged her back. "My dear
        cousin," he pleaded, "do stop a bit! Let me tell you just one thing;}
}
{
}

\Verse{133}
{
    \zA{你的话我都知道了。”}
}
{
    \eA{after that, you may go." "What can you have to tell me?"}
}
{
}

\Verse{134}
{
    \zA{口里说着，却头也不回，竟去了。}
}
{
    \eA{exclaimed Lin Tai-yü, who while wiping her tears, extricated her hand from
        his grasp.}
}
{
}

\Verse{135}
{
    \zA{宝玉望着，只管发起呆来。}
}
{
    \eA{"I know." she cried, "all you have to say." As she spoke, she went away,}
}
{
}

\Verse{136}
{
    \zA{原来方才出来忙了，不曾带得扇子，袭人怕他热， 忙拿了扇子赶来送给他，猛抬头看见黛玉和他站着。}
}
{
    \eA{without even turning her head to cast a glance behind her. As Pao-yü gazed
        at her receding figure, he fell into abstraction. He had, in fact, quitted
        his apartments a few moments back in such precipitate hurry that he had
        omitted to take a fan with him: and Hsi Jen, fearing lest he might suffer
        from the heat,}
}
{
}

\Verse{137}
{
    \zA{一时黛玉走了，他还站着不动， 因而赶上来说道：“你也不带了扇子去，亏了我看见，赶着送来。”}
}
{
    \eA{promptly seized one and ran to find him and give it to him. But upon
        casually raising her head, she espied Lin Tai-yü standing with him. After a
        time, Tai-yü walked away; and as he still remained where he was without
        budging,}
}
{
}

\Verse{138}
{
    \zA{宝玉正出了神， 见袭人和他说话，并未看出是谁，只管呆着脸说道：“好妹妹，我的这个心，从来
        不敢说，今日胆大说出来，就是死了也是甘心的!我为你也弄了一身的病，又不敢 告诉人，只好捱着。}
}
{
    \eA{she approached him. "You left," she said, "without even taking a fan with
        you. Happily I noticed it, and so hurried to catch you up and bring it to
        you." But Pao-yü was so lost in thought that as soon as he caught Hsi Jen's
        voice, he made a dash and clasped her in his embrace, without so much as
        trying to make sure who she was.}
}
{
}

\Verse{139}
{
    \zA{等你的病好了，只怕我的病才得好呢。}
}
{
    \eA{"My dear cousin," he cried, "I couldn't hitherto muster enough courage to
        disclose the secrets of my heart;}
}
{
}

\Verse{140}
{
    \zA{睡里梦里也忘不了你！”}
}
{
    \eA{but on this occasion I shall make bold and give utterance to them.}
}
{
}

\Verse{141}
{
    \zA{袭人听了，惊疑不止，又是怕，又是急，又是臊，连忙推他道：“这是那里的话？}
}
{
    \eA{For you I'm quite ready to even pay the penalty of death. I have too for
        your sake brought ailments upon my whole frame. It's in here! But I haven't
        ventured to breathe it to any one. My only alternative has been to bear it
        patiently,}
}
{
}

\Verse{142}
{
    \zA{你是怎么着了？}
}
{
    \eA{in the hope that when you got all right,}
}
{
}

\Verse{143}
{
    \zA{还不快去吗？”}
}
{
    \eA{I might then perchance also recover.}
}
{
}

\Verse{144}
{
    \zA{宝玉一时醒过来，方知是袭人。}
}
{
    \eA{But whether I sleep, or whether I dream, I never, never forget you." These
        declarations quite dumfoundered Hsi Jen.}
}
{
}

\Verse{145}
{
    \zA{虽然羞的满面紫涨， 却仍是呆呆的，接了扇子，一句话也没有，竟自走去。}
}
{
    \eA{She gave way to incessant apprehensions. All she could do was to shout out:
        "Oh spirits, oh heaven, oh Buddha, he's compassing my death!" Then pushing
        him away from her,}
}
{
}

\Verse{146}
{
    \zA{这里袭人见他去后，想他方才之言必是因黛玉而起，如此看来，倒怕将来难免 不才之事，令人可惊可畏。}
}
{
    \eA{"what is it you're saying?" she asked. "May it be that you are possessed by
        some evil spirit! Don't you quick get yourself off?" This brought Pao-yü to
        his senses at once. He then became aware that it was Hsi Jen,}
}
{
}

\Verse{147}
{
    \zA{却是如何处治，方能免此丑祸？}
}
{
    \eA{and that she had come to bring him a fan. Pao-yü was overpowered with shame;}
}
{
}

\Verse{148}
{
    \zA{想到此间，也不觉呆呆 的发起怔来。}
}
{
    \eA{his whole face was suffused with scarlet; and, snatching the fan out of her
        hands, he bolted away with rapid stride.}
}
{
}

\Verse{149}
{
    \zA{谁知宝钗恰从那边走来，笑道：“大毒日头地下，出什么神呢？”}
}
{
    \eA{When Hsi Jen meanwhile saw Pao-yü effect his escape, "Lin Tai-yü," she
        pondered, "must surely be at the bottom of all he said just now. But from
        what one can see,}
}
{
}

\Verse{150}
{
    \zA{袭 人见问，忙笑说道：“我才见两个雀儿打架，倒很有个玩意儿，就看住了。”}
}
{
    \eA{it will be difficult, in the future, to obviate the occurrence of some
        unpleasant mishap. It's sufficient to fill one with fear and trembling!" At
        this point in her cogitations,}
}
{
}

\Verse{151}
{
    \zA{宝钗 道：“宝兄弟才穿了衣服，忙忙的那里去了？}
}
{
    \eA{she involuntarily melted into tears, so agitated was she; while she secretly
        exercised her mind how best to act so as to prevent this dreadful calamity.}
}
{
}

\Verse{152}
{
    \zA{我要叫住问他呢，只是他慌慌张张的 走过去，竟像没理会我的，所以没问。”}
}
{
    \eA{But while she was lost in this maze of surmises and doubts, Pao-ch'ai
        unexpectedly appeared from the off side. "What!" she smilingly exclaimed,
        "are you dreaming away in a hot broiling sun like this?"}
}
{
}

\Verse{153}
{
    \zA{袭人道：“老爷叫他出去的。”}
}
{
    \eA{Hsi Jen, at this question, hastily returned her smiles. "Those two birds,"
        she answered,}
}
{
}

\Verse{154}
{
    \zA{宝钗听了， 忙说道：“嗳哟，这么大热的天，叫他做什么？}
}
{
    \eA{"were having a fight, and such fun was it that I stopped to watch them."
        "Where is cousin Pao off to now in such a hurry, got up in that fine
        attire?"}
}
{
}

\Verse{155}
{
    \zA{别是想起什么来生了气，叫他出去 教训一场罢？”}
}
{
    \eA{asked Pao-ch'ai, "I just caught sight of him, as he went by. I meant to have
        called out and stopped him, but as he, of late, talks greater rubbish than
        ever,}
}
{
}

\Verse{156}
{
    \zA{袭人笑道：“不是这个，想必有客要会。”}
}
{
    \eA{I didn't challenge him, but let him go past." "Our master," rejoined Hsi
        Jen, "sent for him to go out."}
}
{
}

\Verse{157}
{
    \zA{宝钗笑道：“这个客也 没意思，这么热天不在家里凉快，跑什么！”}
}
{
    \eA{"Ai-yah!" hastily exclaimed Pao-ch'ai, as soon as this remark reached her
        ears. "What does he want him for, on a scalding day like this? Might he not
        have thought of something and got so angry about it as to send for him to
        give him a lecture!"}
}
{
}

\Verse{158}
{
    \zA{袭人笑道：“你可说么！”}
}
{
    \eA{"If it isn't this," added Hsi Jen laughing, "some visitor must, I presume,}
}
{
}

\Verse{159}
{
    \zA{宝钗因问：“云丫头在你们家做什么呢？”}
}
{
    \eA{have come and he wishes him to meet him." "With weather like this," smiled
        Pao-ch'ai, "even visitors afford no amusement!}
}
{
}

\Verse{160}
{
    \zA{袭人笑道：“才说了会子闲话儿， 又瞧了会子我前日粘的鞋帮子，明日还求他做去呢。”}
}
{
    \eA{Why don't they, while this fiery temperature lasts, stay at home, where it's
        much cooler, instead of gadding about all over the place?" "Could you tell
        them so?" smiled Hsi Jen. "What was that girl Hsiang-yün doing in your
        quarters?"}
}
{
}

\Verse{161}
{
    \zA{宝钗听见这话，便两边回头， 看无人来往，笑道：“你这么个明白人，怎么一时半刻的就不会体谅人？}
}
{
    \eA{Pao-ch'ai then asked. "She only came to chat with us on irrelevant matters."
        Hsi Jen replied smiling. "But did you see the pair of shoes I was pasting
        the other day? Well, I meant to ask her to-morrow to finish them for me."
        Pao-chai, at these words,}
}
{
}

\Verse{162}
{
    \zA{我近来看 着云姑娘的神情儿，风里言风里语的听起来，在家里一点儿做不得主。}
}
{
    \eA{turned her head round, first on this side, and then on the other. Seeing
        that there was no one coming or going: "How is it," she smiled, "that you,
        who have so much gumption, don't ever show any respect for people's
        feelings?}
}
{
}

\Verse{163}
{
    \zA{他们家嫌费 用大，竟不用那些针线上的人，差不多儿的东西都是他们娘儿们动手。}
}
{
    \eA{I've been of late keeping an eye on Miss Yün's manner, and, from what I can
        glean from the various rumours afloat, she can't be, in the slightest
        degree, her own mistress at home! In that family of theirs,}
}
{
}

\Verse{164}
{
    \zA{为什么这几 次他来了，他和我说话儿，见没人在跟前，他就说家里累的慌？}
}
{
    \eA{so little can they stand the burden of any heavy expenses that they don't
        employ any needlework-people, and ordinary everyday things are mostly
        attended to by their ladies themselves. (If not), why is it that every time
        she has come to us on a visit, and she and I have had a chat,}
}
{
}

\Verse{165}
{
    \zA{我再问他两句家常 过日子的话，他就连眼圈儿都红了，嘴里含含糊糊待说不说的。}
}
{
    \eA{she at once broached the subject of their being in great difficulties at
        home, the moment she perceived that there was no one present? Yet, whenever
        I went on to ask her a few questions about their usual way of living,}
}
{
}

\Verse{166}
{
    \zA{看他的形景儿，自 然从小儿没了父母是苦的。}
}
{
    \eA{her very eyes grew red, while she made some indistinct reply; but as for
        speaking out, she wouldn't. But when I consider the circumstances in which
        she is placed,}
}
{
}

\Verse{167}
{
    \zA{我看见他也不觉的伤起心来。”}
}
{
    \eA{for she has certainly had the misfortune of being left, from her very
        infancy, without father and mother,}
}
{
}

\Verse{168}
{
    \zA{袭人见说这话，将手一 拍道：“是了。}
}
{
    \eA{the very sight of her is too much for me, and my heart begins to bleed
        within me." "Quite so! Quite so!" observed Hsi Jen,}
}
{
}

\Verse{169}
{
    \zA{怪道上月我求他打十根蝴蝶儿结子，过了那些日子才打发人送来， 还说：‘这是粗打的，且在别处将就使罢；要匀净的，等明日来住着再好生打。’}
}
{
    \eA{clapping her hands, after listening to her throughout. "It isn't strange
        then if she let me have the ten butterfly knots I asked her to tie for me
        only after ever so many days, and if she said that they were coarsely done,
        but that I should make the best of them and use them elsewhere, and that if
        I wanted any nice ones, I should wait until by and bye when she came to stay
        here,}
}
{
}

\Verse{170}
{
    \zA{如今听姑娘这话，想来我们求他，他不好推辞，不知他在家里怎么三更半夜的做呢! 可是我也糊涂了，早知道是这么着，我也不该求他！”}
}
{
    \eA{when she would work some neatly for me. What you've told me now reminds me
        that, as she had found it difficult to find an excuse when we appealed to
        her, she must have had to slave away, who knows how much, till the third
        watch in the middle of the night. What a stupid thing I was!}
}
{
}

\Verse{171}
{
    \zA{宝钗道：“上次他告诉我， 说在家里做活做到三更天，要是替别人做一点半点儿，那些奶奶太太们还不受用 呢。”}
}
{
    \eA{Had I known this sooner, I would never have told her a word about it." "Last
        time;" continued Pao-ch'ai, "she told me that when she was at home she had
        ample to do, that she kept busy as late as the third watch, and that, if she
        did the slightest stitch of work for any other people,}
}
{
}

\Verse{172}
{
    \zA{袭人道：“偏我们那个牛心的小爷，凭着小的大的活计，一概不要家里这些 活计上的人做，我又弄不开这些。”}
}
{
    \eA{the various ladies, belonging to her family, did not like it." "But as it
        happens," explained Hsi Jen, "that mulish-minded and perverse-tempered young
        master of ours won't allow the least bit of needlework, no matter whether
        small or large, to be made by those persons employed to do sewing in the
        household.}
}
{
}

\Verse{173}
{
    \zA{宝钗笑道：“你理他呢!只管叫人做去就是了。”}
}
{
    \eA{And as for me, I have no time to turn my attention to all these things."
        "Why mind him?" laughed Pao-ch'ai. "Simply ask some one to do the work and
        finish."}
}
{
}

\Verse{174}
{
    \zA{袭人道：“那里哄的过他？}
}
{
    \eA{"How could one bamboozle him?" resumed Hsi Jen. "Why, he'll promptly find
        out everything.}
}
{
}

\Verse{175}
{
    \zA{他才是认得出来呢。}
}
{
    \eA{Such a thing can't even be suggested.}
}
{
}

\Verse{176}
{
    \zA{说不得我只好慢慢的累去罢了。”}
}
{
    \eA{The only thing I can do is to quietly slave away, that's all." "You
        shouldn't work so hard,"}
}
{
}

\Verse{177}
{
    \zA{宝钗笑道：“你不必忙，我替你做些就是了。”}
}
{
    \eA{smiled Pao-ch'ai. "What do you say to my doing a few things for you?"}
}
{
}

\Verse{178}
{
    \zA{袭人笑道：“当真的？}
}
{
    \eA{"Are you in real earnest!" ventured Hsi Jen smiling.}
}
{
}

\Verse{179}
{
    \zA{这可就是我 的造化了!晚上我亲自过来――” 一句话未了，忽见一个老婆子忙忙走来，说道：“这是那里说起!金钏儿姑娘 好好儿的投井死了！”}
}
{
    \eA{"Well, in that case, it is indeed a piece of good fortune for me! I'll come
        over myself in the evening." But before she could conclude her reply, she of
        a sudden noticed an old matron come up to her with precipitate step. "Where
        does the report come from," she interposed, "that Miss Chin Ch'uan-erh has
        gone,}
}
{
}

\Verse{180}
{
    \zA{袭人听得，唬了一跳，忙问：“那个金钏儿？”}
}
{
    \eA{for no rhyme or reason, and committed suicide by jumping into the well?"
        This bit of news startled Hsi Jen.}
}
{
}

\Verse{181}
{
    \zA{那老婆子道： “那里还有两个金钏儿呢？}
}
{
    \eA{"Which Chin Ch'uan-erh is it," she speedily inquired. "Where are two Chin
        Ch'uan-erhs to be found!"}
}
{
}

\Verse{182}
{
    \zA{就是太太屋里的。}
}
{
    \eA{rejoined the old matron. "It's the one in our Mistress,'}
}
{
}

\Verse{183}
{
    \zA{前日不知为什么撵出去，在家里哭天 抹泪的，也都不理会他，谁知找不着他，才有打水的人说那东南角上井里打水，见
        一个尸首，赶着叫人打捞起来，谁知是他!他们还只管乱着要救，那里中用了呢？”}
}
{
    \eA{Madame Wang's, apartments, who was the other day sent away for something or
        other, I don't know what. On her return home, she raised her groans to the
        skies and shed profuse tears, but none of them worried their minds about
        her,}
}
{
}

\Verse{184}
{
    \zA{宝钗道：“这也奇了！”}
}
{
    \eA{until, who'd have thought it, they could see nothing of her.}
}
{
}

\Verse{185}
{
    \zA{袭人听说，点头赞叹，想素日同气之情，不觉流下泪来。}
}
{
    \eA{A servant, however, went just now to draw water and he says that 'while he
        was getting it from the well in the south-east corner, he caught sight of a
        dead body, that he hurriedly called men to his help, and that when they
        fished it out,}
}
{
}

\Verse{186}
{
    \zA{宝钗听见这话，忙向王夫人处来安慰。}
}
{
    \eA{they unexpectedly found that it was she, but that though they bustled about
        trying to bring her round, everything proved of no avail'" "This is odd!"}
}
{
}

\Verse{187}
{
    \zA{这里袭人自回去了。}
}
{
    \eA{Pao-ch'ai exclaimed. The moment Hsi Jen heard the tidings, she shook her
        head and moaned.}
}
{
}

\Verse{188}
{
    \zA{宝钗来至王夫人房里，只见鸦雀无闻，独有王夫人在里间房内坐着垂泪。}
}
{
    \eA{At the remembrance of the friendship, which had ever existed between them,
        tears suddenly trickled down her cheeks. And as for Pao-ch'ai, she listened
        to the account of the accident and then hastened to Madame Wang's quarters
        to try and afford her consolation.}
}
{
}

\Verse{189}
{
    \zA{宝钗 便不好提这事，只得一旁坐下。}
}
{
    \eA{Hsi Jen, during this interval, returned to her room. But we will leave her
        without further notice,}
}
{
}

\Verse{190}
{
    \zA{王夫人便问：“你打那里来？”}
}
{
    \eA{and explain that when Pao-ch'ai reached the interior of Madame Wang's home,}
}
{
}

\Verse{191}
{
    \zA{宝钗道：“打园里 来。”}
}
{
    \eA{she found everything plunged in perfect stillness.}
}
{
}

\Verse{192}
{
    \zA{王夫人道：“你打园里来，可曾见你宝兄弟？”}
}
{
    \eA{Madame Wang was seated all alone in the inner chamber indulging her sorrow.}
}
{
}

\Verse{193}
{
    \zA{宝钗道：“才倒看见他了： 穿着衣裳出去了，不知那里去。”}
}
{
    \eA{But such difficulties did Pao-ch'ai experience to allude to the occurrence,
        that her only alternative was to take a seat next to her. "Where do you come
        from?"}
}
{
}

\Verse{194}
{
    \zA{王夫人点头叹道：“你可知道一件奇事？}
}
{
    \eA{asked Madame Wang. "I come from inside the garden," answered Pao-ch'ai. "As
        you come from the garden,"}
}
{
}

\Verse{195}
{
    \zA{金钏儿 忽然投井死了！”}
}
{
    \eA{Madame Wang inquired, "did you see anything of your cousin Pao-yü?"}
}
{
}

\Verse{196}
{
    \zA{宝钗见说，道：“怎么好好儿的投井？}
}
{
    \eA{"I saw him just now," Pao-ch'ai replied, "go out, dressed up in his
        fineries.}
}
{
}

\Verse{197}
{
    \zA{这也奇了。”}
}
{
    \eA{But where he is gone to,}
}
{
}

\Verse{198}
{
    \zA{王夫人道： “原是前日他把我一件东西弄坏了，我一时生气，打了他两下子，撵了下去。}
}
{
    \eA{I don't know." "Have you perchance heard of any strange occurrence?" asked
        Madame Wang, while she nodded her head and sighed. "Why, Chin Ch'uan Erh
        jumped into the well and committed suicide."}
}
{
}

\Verse{199}
{
    \zA{我只 说气他几天，还叫他上来，谁知他这么气性大，就投井死了。}
}
{
    \eA{"How is it that she jumped into the well when there was nothing to make her
        do so?" Pao-ch'ai inquired. "This is indeed a remarkable thing!" "The fact
        is," proceeded Madame Wang,}
}
{
}

\Verse{200}
{
    \zA{岂不是我的罪过！”}
}
{
    \eA{"that she spoilt something the other day, and in a sudden fit of temper,}
}
{
}

\Verse{201}
{
    \zA{宝钗笑道：“姨娘是慈善人，固然是这么想。}
}
{
    \eA{I gave her a slap and sent her away, simply meaning to be angry with her for
        a few days and then bring her in again.}
}
{
}

\Verse{202}
{
    \zA{据我看来，他并不是赌气投井，多半 他下去住着，或是在井傍边儿玩，失了脚掉下去的。}
}
{
    \eA{But, who could have ever imagined that she had such a resentful temperament
        as to go and drown herself in a well! And is not this all my fault?" "It's
        because you are such a kind-hearted person, aunt," smiled Pao-ch'ai,}
}
{
}

\Verse{203}
{
    \zA{他在上头拘束惯了，这一出去 自然要到各处去玩玩逛逛儿，岂有这样大气的理？}
}
{
    \eA{"that such ideas cross your mind! But she didn't jump into the well when she
        was in a tantrum; so what must have made her do so was that she had to go
        and live in the lower quarters. Or, she might have been standing in front of
        the well,}
}
{
}

\Verse{204}
{
    \zA{纵然有这样大气，也不过是个糊 涂人，也不为可惜。”}
}
{
    \eA{and her foot slipped, and she fell into it. While in the upper rooms, she
        used to be kept under restraint, so when this time she found herself
        outside,}
}
{
}

\Verse{205}
{
    \zA{王夫人点头叹道：“虽然如此，到底我心里不安！”}
}
{
    \eA{she must, of course, have felt the wish to go strolling all over the place
        in search of fun. How could she have ever had such a fiery disposition?}
}
{
}

\Verse{206}
{
    \zA{宝钗笑 道：“姨娘也不劳关心。}
}
{
    \eA{But even admitting that she had such a temper, she was, after all, a stupid
        girl to do as she did;}
}
{
}

\Verse{207}
{
    \zA{十分过不去，不过多赏他几两银子发送他，也就尽了主仆 之情了。”}
}
{
    \eA{and she doesn't deserve any pity." "In spite of what you say," sighed Madame
        Wang, shaking her head to and fro, "I really feel unhappy at heart."}
}
{
}

\Verse{208}
{
    \zA{王夫人道：“才刚我赏了五十两银子给他妈，原要还把你姐妹们的新衣 裳给他两件装裹，谁知可巧都没有什么新做的衣裳，只有你林妹妹做生日的两套。}
}
{
    \eA{"You shouldn't, aunt, distress your mind about it!" Pao-ch'ai smiled. "Yet,
        if you feel very much exercised, just give her a few more taels than you
        would otherwise have done, and let her be buried. You'll thus carry out to
        the full the feelings of a mistress towards her servant." "I just now gave
        them fifty taels for her,"}
}
{
}

\Verse{209}
{
    \zA{我想你林妹妹那孩子，素日是个有心的，况且他也三灾八难的，既说了给他作生日， 这会子又给人去装裹，岂不忌讳？}
}
{
    \eA{pursued Madame Wang. "I also meant to let them have some of your cousin's
        new clothes to enshroud her in. But, who'd have thought it, none of the
        girls had, strange coincidence, any newly-made articles of clothing; and
        there were only that couple of birthday suits of your cousin Lin's.}
}
{
}

\Verse{210}
{
    \zA{因这么着，我才现叫裁缝赶着做一套给他。}
}
{
    \eA{But as your cousin Lin has ever been such a sensitive child and has always
        too suffered and ailed,}
}
{
}

\Verse{211}
{
    \zA{要是 别的丫头，赏他几两银子，也就完了。}
}
{
    \eA{I thought it would be unpropitious for her, if her clothes were also now
        handed to people to wrap their dead in,}
}
{
}

\Verse{212}
{
    \zA{金钏儿虽然是个丫头，素日在我跟前，比我 的女孩儿差不多儿！”}
}
{
    \eA{after she had been told that they were given her for her birthday. So I
        ordered a tailor to get a suit for her as soon as possible. Had it been any
        other servant-girl,}
}
{
}

\Verse{213}
{
    \zA{口里说着，不觉流下泪来。}
}
{
    \eA{I could have given her a few taels and have finished.}
}
{
}

\Verse{214}
{
    \zA{宝钗忙道：“姨娘这会子何用叫 裁缝赶去。}
}
{
    \eA{But Chin Ch'uan-erh was, albeit a servant-maid, nearly as dear to me as if
        she had been a daughter of mine."}
}
{
}

\Verse{215}
{
    \zA{我前日倒做了两套，拿来给他，岂不省事？}
}
{
    \eA{Saying this, tears unwittingly ran down from her eyes. "Aunt!" vehemently
        exclaimed Pao-ch'ai.}
}
{
}

\Verse{216}
{
    \zA{况且他活的时候儿也穿过我 的旧衣裳，身量也相对。”}
}
{
    \eA{"What earthly use is it of hurrying a tailor just now to prepare clothes for
        her? I have a couple of suits I made the other day and won't it save trouble
        were I to go and bring them for her?}
}
{
}

\Verse{217}
{
    \zA{王夫人道：“虽然这样，难道你不忌讳？”}
}
{
    \eA{Besides, when she was alive, she used to wear my old clothes. And what's
        more our figures are much alike."}
}
{
}

\Verse{218}
{
    \zA{宝钗笑道： “姨娘放心，我从来不计较这些。”}
}
{
    \eA{"What you say is all very well," rejoined Madame Wang; "but can it be that
        it isn't distasteful to you?"}
}
{
}

\Verse{219}
{
    \zA{一面说，一面起身就走。}
}
{
    \eA{"Compose your mind," urged Pao-ch'ai with a smile.}
}
{
}

\Verse{220}
{
    \zA{王夫人忙叫了两个人 跟宝钗去。}
}
{
    \eA{"I have never paid any heed to such things." As she spoke, she rose to her
        feet and walked away.}
}
{
}

\Verse{221}
{
    \zA{一时宝钗取了衣服回来，只见宝玉在王夫人旁边坐着垂泪。}
}
{
    \eA{Madame Wang then promptly called two servants. "Go and accompany Miss Pao!"
        she said. In a brief space of time, Pao-ch'ai came back with the clothes,}
}
{
}

\Verse{222}
{
    \zA{王夫人正才说他， 因见宝钗来了，就掩住口不说了。}
}
{
    \eA{and discovered Pao-yü seated next to Madame Wang, all melted in tears.
        Madame Wang was reasoning with him. At the sight of Pao-ch'ai,}
}
{
}

\Verse{223}
{
    \zA{宝钗见此景况，察言观色，早知觉了七八分。}
}
{
    \eA{she, at once, desisted. When Pao-ch'ai saw them go on in this way, and came
        to weigh their conversation and to scan the expression on their
        countenances,}
}
{
}

\Verse{224}
{
    \zA{于 是将衣服交明王夫人，王夫人便将金钏儿的母亲叫来拿了去了。}
}
{
    \eA{she immediately got a pretty correct insight into their feelings. But
        presently she handed over the clothes, and Madame Wang sent for Chin
        Ch'uan-erh's mother,}
}
{
}

\Verse{225}
{
    \zA{后事如何，下回分解。}
}
{
    \eA{to take them away. But, reader,}
}
{
}

\Verse{226}
{
    \zA{(本章完)}
}
{
    \eA{you will have to peruse the next chapter for further details.}
}
{
}
